% GENERATED FILE. Edit canonical lessons, then run npm run generate:books.
% canonical-source-hash: fnv1a64:4be32a5ba4708f44
% canonical-lessons: PT-C436-indicacoes, PT-C436-rotunda, PT-C436-contornar, PT-C436-entroncamento, PT-C436-prosseguir

\chapter{Pedir indicações}
\label{ch:pt-pedir-indicacoes}
\hlchaptermodality{436}

\begin{chapteropening}
\textbf{By the end of this chapter:} \emph{I can ask for and follow directions: directions, the roundabout, going round, the road junction and carrying on.}

Complete the last lesson of chapter 436: I can ask for and follow directions: directions, the roundabout, going round, the road junction and carrying on.
\end{chapteropening}

\section[as indicações]{as indicações — the directions; the instructions}

\label{lesson:PT-C436-indicacoes}



\begin{quote}
Before the new one: say the Portuguese for the travel plan, then the Portuguese for to locate.
\end{quote}

\subsection*{You'll want to know: as indicações}
\textbf{as indicações} — \textquotedblleft{}the directions; the instructions\textquotedblright{}.

\begin{quote}
Um senhor deu-nos indicações para a estação.
\end{quote}

\begin{quote}
\emph{A man gave us directions to the station.}
\end{quote}

Singular \emph{a indicação} — a sign, a hint.

\begin{cousinweb}
\emph{indicação} is a learned borrowing from Latin \emph{indicātiō} 'pointing out' (EtymDB), built on \emph{indicāre} 'to point out'.
\end{cousinweb}

\begin{grammarlens}[title={What you've built}]
One more word for telling, explaining and describing at B1.
\end{grammarlens}

\subsection*{Your turn}
\begin{itemize}
\raggedright
  \item \emph{Say it:} \emph{as indicações}
  \item \emph{Say it:} \emph{as indicações}, once more
  \item \emph{Say it:} a sentence of your own with \emph{as indicações}
  \item \emph{Recall:} say the Portuguese for the travel plan, then the Portuguese for to locate, then say \emph{as indicações} again
\end{itemize}

\begin{tcolorbox}[breakable,colback=teal!4,colframe=teal!35!black,title={Before you move on}]
What does as indicações mean? (\textquotedblleft{}The directions; the instructions\textquotedblright{}.) Say it once more.
\end{tcolorbox}
\section[a rotunda]{a rotunda — the roundabout}

\label{lesson:PT-C436-rotunda}



\begin{quote}
Before the new one: say the Portuguese for to locate, then the Portuguese for the directions.
\end{quote}

\subsection*{You'll want to know: a rotunda}
\textbf{a rotunda} — \textquotedblleft{}the roundabout\textquotedblright{}.

\begin{quote}
Na rotunda, apanhamos a segunda saída.
\end{quote}

\begin{quote}
\emph{At the roundabout, we take the second exit.}
\end{quote}

In Brazil: \emph{a rotatória}.

\begin{cousinweb}
\emph{rotunda} is a learned borrowing from Latin \emph{rotunda} 'round' (EtymDB has \emph{rotundus}, from \emph{rota} 'wheel'); the inherited form of the same word is \emph{redondo}.
\end{cousinweb}

\begin{grammarlens}[title={What you've built}]
One more word for telling, explaining and describing at B1.
\end{grammarlens}

\subsection*{Your turn}
\begin{itemize}
\raggedright
  \item \emph{Say it:} \emph{a rotunda}
  \item \emph{Say it:} \emph{a rotunda}, once more
  \item \emph{Say it:} a sentence of your own with \emph{a rotunda}
  \item \emph{Recall:} say the Portuguese for to locate, then the Portuguese for the directions, then say \emph{a rotunda} again
\end{itemize}

\begin{tcolorbox}[breakable,colback=teal!4,colframe=teal!35!black,title={Before you move on}]
What does a rotunda mean? (\textquotedblleft{}The roundabout\textquotedblright{}.) Say it once more.
\end{tcolorbox}
\section[contornar]{contornar — to go round; to get round (a problem)}

\label{lesson:PT-C436-contornar}



\begin{quote}
Before the new one: say the Portuguese for the directions, then the Portuguese for the roundabout.
\end{quote}

\subsection*{You'll want to know: contornar}
\textbf{contornar} — \textquotedblleft{}to go round; to get round (a problem)\textquotedblright{}.

\begin{quote}
Contornamos a praça e chegamos ao museu.
\end{quote}

\begin{quote}
\emph{We go round the square and reach the museum.}
\end{quote}

\emph{Contornar um problema} — to get round a problem.

\begin{cousinweb}
\emph{contornar} is built from \emph{con-} and \emph{tornar} 'to turn', which is inherited from Latin \emph{tornāre} 'to turn on a lathe' (EtymDB).
\end{cousinweb}

\begin{grammarlens}[title={What you've built}]
One more word for telling, explaining and describing at B1.
\end{grammarlens}

\subsection*{Your turn}
\begin{itemize}
\raggedright
  \item \emph{Say it:} \emph{contornar}
  \item \emph{Say it:} \emph{contornar}, once more
  \item \emph{Say it:} a sentence of your own with \emph{contornar}
  \item \emph{Recall:} say the Portuguese for the directions, then the Portuguese for the roundabout, then say \emph{contornar} again
\end{itemize}

\begin{tcolorbox}[breakable,colback=teal!4,colframe=teal!35!black,title={Before you move on}]
What does contornar mean? (\textquotedblleft{}To go round; to get round (a problem)\textquotedblright{}.) Say it once more.
\end{tcolorbox}
\section[o entroncamento]{o entroncamento — the road junction; the railway junction}

\label{lesson:PT-C436-entroncamento}



\begin{quote}
Before the new one: say the Portuguese for the roundabout, then the Portuguese for to go round.
\end{quote}

\subsection*{You'll want to know: o entroncamento}
\textbf{o entroncamento} — \textquotedblleft{}the road junction; the railway junction\textquotedblright{}.

\begin{quote}
No entroncamento, viramos à esquerda.
\end{quote}

\begin{quote}
\emph{At the junction, we turn left.}
\end{quote}

Where two roads or railway lines meet like branches on a trunk.

\begin{cousinweb}
\emph{entroncamento} is built on \emph{entroncar} 'to join together', on \emph{tronco} 'trunk', which is inherited from Latin \emph{truncus} (EtymDB).
\end{cousinweb}

\begin{grammarlens}[title={What you've built}]
One more word for telling, explaining and describing at B1.
\end{grammarlens}

\subsection*{Your turn}
\begin{itemize}
\raggedright
  \item \emph{Say it:} \emph{o entroncamento}
  \item \emph{Say it:} \emph{o entroncamento}, once more
  \item \emph{Say it:} a sentence of your own with \emph{o entroncamento}
  \item \emph{Recall:} say the Portuguese for the roundabout, then the Portuguese for to go round, then say \emph{o entroncamento} again
\end{itemize}

\begin{tcolorbox}[breakable,colback=teal!4,colframe=teal!35!black,title={Before you move on}]
What does o entroncamento mean? (\textquotedblleft{}The road junction; the railway junction\textquotedblright{}.) Say it once more.
\end{tcolorbox}
\section[prosseguir]{prosseguir — to carry on, to continue}

\label{lesson:PT-C436-prosseguir}



\begin{quote}
Before the new one: say the Portuguese for to go round, then the Portuguese for the road junction.
\end{quote}

\subsection*{You'll want to know: prosseguir}
\textbf{prosseguir} — \textquotedblleft{}to carry on, to continue\textquotedblright{}.

\begin{quote}
Depois da ponte, prosseguimos em frente.
\end{quote}

\begin{quote}
\emph{After the bridge, we carried on straight ahead.}
\end{quote}

A little formal; people also say \emph{seguir}.

\begin{cousinweb}
\emph{prosseguir} is a learned borrowing from Latin \emph{prōsequī} 'to follow on' (EtymDB), reshaped on \emph{seguir} 'to follow'.
\end{cousinweb}

\begin{grammarlens}[title={What you've built}]
One more word for telling, explaining and describing at B1.
\end{grammarlens}

\subsection*{Your turn}
\begin{itemize}
\raggedright
  \item \emph{Say it:} \emph{prosseguir}
  \item \emph{Say it:} \emph{prosseguir}, once more
  \item \emph{Say it:} a sentence of your own with \emph{prosseguir}
  \item \emph{Recall:} say the Portuguese for to go round, then the Portuguese for the road junction, then say \emph{prosseguir} again
\end{itemize}

\begin{tcolorbox}[breakable,colback=teal!4,colframe=teal!35!black,title={Before you move on}]
What does prosseguir mean? (\textquotedblleft{}To carry on, to continue\textquotedblright{}.) Say it once more.
\end{tcolorbox}
